% Source: https://artofproblemsolving.com/community/c3h2404259_i_need_help_may_anyone_help_me
% Stable identity: aops-topic-2404259-q01
Write $n=a^2-b^2=(a-b)(a+b)=de$ with $d=a-b$ and $e=a+b$. If $a,b$ are positive integers, then $1\le d<e$ and $d,e$ have the same parity (their sum $2a$ is even). Conversely, if $n=de$ with $1\le d<e$ and $d\equiv e\pmod 2$, then $a=\dfrac{e+d}{2}$ and $b=\dfrac{e-d}{2}$ are positive integers with $a^2-b^2=n$. So $n$ is representable exactly when it has such a factorisation.
\begin{itemize}
\item $n\equiv2\pmod4$: two odd factors give an odd product and two even factors give a multiple of $4$, so $n$ is never representable. These are $n=2,6,10,\dots,98$: $25$ numbers.
\item $n$ odd: for $n\ge3$ take $d=1$, $e=n$. For $n=1$ the only factorisation is $1\cdot1$, with $d=e$ (it gives $b=0$), so $n=1$ is not representable.
\item $4\mid n$, $n=4m$: for $m\ge2$ take $d=2$, $e=2m>2$. For $n=4$ the factorisations are $1\cdot4$ (different parity) and $2\cdot2$ ($d=e$), so $n=4$ is not representable.
\end{itemize}
The numbers that cannot be written in this way are the $25$ numbers $\equiv2\pmod4$ together with $1$ and $4$: $\boxed{27}$ numbers.
